\documentclass[10pt,a4paper]{amsart}
\usepackage{amsmath,amssymb,amsthm,graphicx}
\usepackage{amsmath}
\usepackage{amssymb}
\usepackage{hyperref}
\renewcommand{\theenumi}{\Alph{enumi}}
\renewcommand{\thesubsection}{\thesection\Alph{subsection}}

%%%%%%%%%%%%%%%%%%%%%%%%%%%%%%%%%%%%%%%%%%
% ENVIRONMENTS
%%%%%%%%%%%%%%%%%%%%%%%%%%%%%%%%%%%%%%%%%%%
\newtheorem{theorem}{Theorem}[section]
\newtheorem{prop}[theorem]{Proposition}
\newtheorem{lemma}[theorem]{Lemma}
\newtheorem{remark}[theorem]{Remark}
\newtheorem{corollary}[theorem]{Corollary}
\newtheorem{definition}[theorem]{Definition}
\newtheorem{example}[theorem]{Example}
%
\def\beq{\begin{equation}}
\def\eeq{\end{equation}}
\def\be{\begin{equation*}}
\def\ee{\end{equation*}}
%%%%%%%%%%%%%%%%%%%%%%%%%
% Beginning of the text %
%%%%%%%%%%%%%%%%%%%%%%%%%
%\begin{document}



%\documentclass[10pt]{article}
\usepackage{amssymb}
\usepackage{amsfonts}
%\topmargin 0.3in
%\textwidth 6in
%\textheight 9in%\baselineskip
%\hoffset=-50pt
%\voffset=-2pt
%\newenvironment{proof}[1][Proof]
%{\ \rule{0.5em}{0.5em}}


\title{Dedekind-finite Cardinals having countable partitions}
\author{Supakun Panasawatwong and J K Truss} 
\date{This paper is an amplified version of part of the first author's PhD thesis at the University of Leeds}%\footnotemark}
\title{Local embeddability versus embeddability under choice, and local isomorphism versus elementary equivalence}
\author{Supakun Panasawatwong, J. K. Truss}
\date{arxiv:2204.05700v1; CC BY 4.0}
\begin{document}
\maketitle
\noindent\textit{Editorial note.} All delivered mathematical text is the author's own, copied from the pinned publisher source sources/dedekind-finite-trees-2204.05700v1/source0.tex; no statement or proof step was rewritten, reordered or removed.\par\medskip
\noindent\textit{Editorial note.} Preamble provenance: the retained preamble is the source's own layout adapter, if any, plus the author's macro block at lines 5-44 of sources/dedekind-finite-trees-2204.05700v1/source0.tex (everything between the publisher class metadata and \textbackslash{}begin\{document\}, comments excluded). Nothing was added, no macro was invented, and the publisher document class is deliberately not reproduced; the excerpt serves only as preamble, never as a span coordinate.\par\medskip
\noindent\textit{Editorial note.} External references are navigation pointers only; every deliberately kept step of the author's argument stays inside the counted spans. The displayed numbers of those pointers are the author's own numbering under the author's \textbackslash{}newtheorem declarations (the source has no subsections, so numbering is section.index), computed mechanically by the lane's scanner.\par\medskip
\noindent\textit{Editorial note.} The counted proof contains no reference to any labelled item and no bibliography citation, so the card needs neither a navigation stub nor a bibliography excerpt.\par\medskip
\noindent\textit{Editorial note.} Source metadata: this directory stages no SOURCE.json; the title and author list were taken from the publisher abstract page kept beside it, and the card points at the per-card licence reuse record generated automatically because the staged licence evidence carries no status field.\par\medskip
\noindent\textit{Allowed standard prerequisites (editorial).}\begin{itemize}
\item Konig's lemma for finitely branching infinite trees. Used inside the counted proof at lines 640-644. Produces an infinite branch in the tree of partial embeddings, whose union is the required embedding.
\item First order sentences capturing the first n levels of an omega-tree up to isomorphism. Used inside the counted proof at lines 648-654. Converts elementary equivalence into local isomorphism by transporting the finite initial segments.
\item The back-and-forth (Ehrenfeucht-Fraisse) argument for elementary equivalence. Used inside the counted proof at lines 656-659. Proves the converse direction, with the player's strategy built from the given local isomorphisms.
\end{itemize}\medskip
\begin{lemma} (i) Assuming the axiom of choice, for any two $\omega$-trees with finite levels, $T_1$ 
is locally embeddable in $T_2$ if and only if it is embeddable in $T_2$.

(ii) Assuming the axiom of choice, $T_1$ and $T_2$ are locally isomorphic if and only if they
are isomorphic.

(iii) $T_1$ and $T_2$ are locally isomorphic if and only if they are elementarily equivalent. \label{4.4}  \end{lemma}

\begin{proof}  
(i) Let $T_1$ be locally embeddable in $T_2$. Let $P_n$ be the family of all embeddings of the union of the
first $n$ levels of $T_1$ into the union of the first $n$ levels of $T_2$. Let $P = \bigcup_{n \in \omega}P_n$
be partially ordered by extension. Then $P$ is a tree with $n$th level $P_n$, each of which is finite, and by
assumption non-empty. By K\"onig's Lemma, $P$ has an infinite branch, and the union of this branch provides an
embedding of $T_1$ into $T_2$.

(ii) This follows by the same proof. 

(iii) Assume first that $T_1$ and $T_2$ are elementarily equivalent. For any $n$ there is a first order sentence
capturing the first $n$ levels of $T_1$ up to isomorphism. More precisely, if $N$ is the number of nodes in these
levels (clearly a finite number), the sentence says that there exist $N$ points, which are ordered in the
correct way, and any other $n$ points which are linearly ordered and form an initial segment are contained in
the $N$ first mentioned. Since $T_1$ and $T_2$ are elementarily equivalent, there must be corresponding $N$ points
of $T_2$, which are ordered in the same way as those in $T_1$, and which form the union of the first $n$ levels.
Hence the unions of the first $n$ levels of $T_1$ and $T_2$ are isomorphic.

Conversely, if $T_1$ and $T_2$ are locally isomorphic, then we can use back-and-forth (an
Ehrenfeucht-Fra\"iss\'e game) to see that they are elementarily equivalent. Suppose that this has $n$
moves. Then player II has a winning strategy, in which he uses given isomorphisms between the unions of
finitely many levels to decide what to do next.          \end{proof}

\end{document}
